\section{The finite generator converges to the realized completed generator}
\label{sec:pggl}

Retain the original increasing remainder frame of
$E=\mathbb C[S]/(\chi_{h,k})$, every primary nilpotent, and
$T_t=A+tR-kI/2$ from PGGT.1, with $t\in\mathbb C$.
PGMB.4--5 and PGRT.6 prove
\[
 \widehat G_N=\Omega_k+T_N,\quad
 T_N=f_N^*f_N=\sum_{j=N}^{\infty}a_j^{-1}\alpha_j^*\alpha_j,
 \quad T_N\longrightarrow0
 \quad(N\to\infty;\ h,k\text{ fixed}).
\tag{PGGL.1}\label{eq:pggl-tail}
\]
The symbol $T_N$ denotes this positive form; the original period
generator is always written $T_t$. Their indices and domains fix
which object is meant. Define the $\Omega_k$-self-adjoint operator
$B_N=\Omega_k^{-1}T_N\succeq_{\Omega_k}0$.
Its full spectral decomposition is $\sum_i b_{N,i}P_{N,i}$ with
$P_{N,i}^*\Omega_k=\Omega_kP_{N,i}$. Transporting the complete
eigenspaces of $\Omega_k^{-1/2}T_N\Omega_k^{-1/2}$ proves this
decomposition in the original frame. Put
\[
 S_N=\sum_i\sqrt{1+b_{N,i}}P_{N,i},\quad
 S_N^{-1}=\sum_i(1+b_{N,i})^{-1/2}P_{N,i},\quad
 S_N^*\Omega_k S_N=\widehat G_N.
\tag{PGGL.2}\label{eq:pggl-isometry}
\]
Multiplication of the orthogonal projections proves the inverse
and Gram identities, including all repeated eigenspaces. Thus
$S_N$ is the specified isometry from $(E,\widehat G_N)$ to
$(E,\Omega_k)$; the original form is not replaced by an identity.

For $G>0$ let $X_G(t)=T_t^{\dagger_G}+T_t$ and
$\epsilon_G(t)=\|X_G(t)\|_G$. Applying the actual isometry gives
the complete error
\[
\begin{gathered}
 S_NX_{\widehat G_N}(t)S_N^{-1}-X_{\Omega_k}(t)
     =Z_N(t)+Z_N(t)^{\dagger_{\Omega_k}},\quad
 Z_N(t)=[S_N,T_t]S_N^{-1},\\
 \|Z_N(t)+Z_N(t)^{\dagger_{\Omega_k}}\|_{\Omega_k}
 \leq\frac{\|[B_N,A]+t[B_N,R]\|_{\Omega_k}}
                       {1+\min_i b_{N,i}}.
\end{gathered}
\tag{PGGL.3}\label{eq:pggl-error}
\]
For the identity, use
$T_t^{\dagger_{\widehat G_N}}=S_N^{-2}T_t^{\dagger_{\Omega_k}}S_N^2$
and expand its two ordered terms. For the bound, the convergent
integral
$\int_0^\infty e^{-uS_N}[B_N,T_t]e^{-uS_N}\,du$
has $(i,j)$ block
$(b_{N,i}-b_{N,j})/(\sqrt{1+b_{N,i}}+\sqrt{1+b_{N,j}})
P_{N,i}T_tP_{N,j}$, hence equals $[S_N,T_t]$.
Its norm is at most the commutator norm divided by
$2\sqrt{1+\min b_{N,i}}$. Multiplication by $S_N^{-1}$ and
addition of the adjoint prove the stated factor one. Both $t$
and $\bar t$ remain in the error through its adjoint.

For every specified finite radius $r\geq0$, the elementary
operator-norm bounds on the two commutators now give
\[
\begin{aligned}
 \sup_{|t|\leq r}|\epsilon_{\widehat G_N}(t)-\epsilon_{\Omega_k}(t)|
 &\leq 2\|B_N\|_{\Omega_k}
                 (\|A\|_{\Omega_k}+r\|R\|_{\Omega_k})\\
 &\leq2(\|A\|_{\Omega_k}+r\|R\|_{\Omega_k})
       \sum_{j=N}^{\infty}
                  \frac{\alpha_j\Omega_k^{-1}\alpha_j^*}{a_j}
 \longrightarrow0.
\end{aligned}
\tag{PGGL.4}\label{eq:pggl-uniform-period-limit}
\]
The first inequality uses the reverse triangle inequality after
PGGL.2, and $1+\min b_{N,i}\geq1$. For the second, a positive
matrix has largest eigenvalue at most its trace. In the original
coordinates this trace is
$\operatorname{tr}(\Omega_k^{-1}T_N)$, whose complete convergent
tail is PGRT.6. The fixed positive matrix $\Omega_k$ and the
fixed operators $A,R$ make the last limit zero. This proves
local uniformity in the original complex period parameter; it
asserts no uniformity in the separate tensor degree $k$.

The receiving operator is the already constructed actual Hilbert
realization, not an unspecified limiting action. Namely let
$\mathfrak U_k=(L_{\rm low},D_{\rm mix})$ with
$\mathfrak U_k^*\mathfrak U_k=\Omega_k$. Then
$\widetilde T_t=\mathfrak U_kT_t\Omega_k^{-1}\mathfrak U_k^*$
intertwines $T_t$ on $\mathfrak U_k(E)$, vanishes on that image's
orthogonal complement, and has kernel
$\mathfrak U_k(\ker T_t)\oplus^\perp\mathfrak U_k(E)^\perp$.
The Gram identity proves these statements by decomposing every
target vector as $\mathfrak U_kx+v_\perp$ with
$x=\Omega_k^{-1}\mathfrak U_k^*v$.
Taking adjoints shows
$\|\widetilde T_t+\widetilde T_t^*\|=\epsilon_{\Omega_k}(t)$.
Consequently PGGL.4 is a quantitative comparison from the finite
original graph generators to this specified mixed-source generator,
with its full kernel, period parameter, and relation-tail error.
Every coefficient map retains its coordinatewise lift to the
original fixed support masks and keeps external absence separate
from each present zero vector.
